\chapter{Introduction}

\section{Background and Motivation}

Grammatical Error Correction (GEC) is a task in Natural Language Processing (NLP) concerned with automatically detecting and correcting errors in learner writing. The field has moved from rule-based systems, to statistical machine translation, and later to neural sequence-to-sequence models \cite{10.1145/3474840}. This shift matters for educational writing support because newer systems can handle a wider range of learner errors, but they also make less transparent editing decisions.

More recently, large language models (LLMs) have been used for zero-shot and few-shot GEC \cite{loem2023exploringeffectivenessgpt3grammatical}. These models can respond to correction prompts without being fine-tuned for the specific dataset. This behaviour is linked to large-scale generative pre-training and in-context learning, which allow the model to produce fluent revisions from instructions and examples \cite{NEURIPS2020_1457c0d6}.

This is useful for rapid experimentation, but it also creates the problem studied in this dissertation: a fluent output is not always a minimal correction. LLMs may improve the surface quality of a sentence while changing wording or structure more than a learner-facing GEC system should.

\section{Problem Statement}

Rather than performing targeted grammatical edits, LLMs often produce broader rewrites that modify lexical choice, tone, or syntactic structure beyond what is strictly necessary. This phenomenon, referred to as \textit{over-correction}, reflects a bias towards producing fluent and natural-sounding text.

Although such revisions may improve readability, they can introduce unnecessary changes that obscure the original structure of learner writing. In the context of second-language (L2) learning, this behaviour is problematic. Effective feedback should not only correct errors but also preserve the learner's original wording, allowing them to identify and reflect on specific linguistic mistakes.

Excessive rewriting reduces transparency and may limit opportunities for error noticing, which is a key component of language learning. Standard GEC metrics, such as ERRANT-based precision, recall, and $F_{0.5}$, are still useful, but they mainly measure agreement with reference corrections and do not fully account for over-correction risk.

For this reason, the project evaluates prompts using both grammatical accuracy and the amount of source-output rewriting.

\section{Aims and Objectives}

The primary aim of this dissertation is to investigate how prompt-based strategies influence potential over-correction in LLM-based grammatical error correction while maintaining correction performance.

To achieve this aim, the study pursues the following objectives:

\begin{itemize}
    \item to analyse the nature of over-correction in LLM-generated outputs;
    \item to develop and evaluate multiple prompting strategies, including instruction-based, role-based, and few-shot approaches;
    \item to compare prompting strategies using both grammatical accuracy and rewriting-based measures;
    \item to introduce and apply OCI to provide an indicator of over-correction risk by combining divergence with fluency and meaning-preservation signals;
    \item to examine the trade-off between correction performance and over-correction;
    \item to investigate whether sentence length influences correction behaviour;
    \item to evaluate the robustness of OCI under alternative weighting schemes.
\end{itemize}

Together, these objectives allow the dissertation to compare prompt designs in terms of both correction performance and over-correction risk.

\section{Scope of the Study}

This study focuses on evaluating prompt engineering within a controlled experimental setting. A single OpenAI GPT-4o configuration is used throughout the experiments, with temperature 0.0, top-p 1.0, a maximum output length of 128 tokens, and seed 1234. Keeping these settings fixed helps ensure that observed differences in output can be attributed mainly to prompt design rather than variation across models or decoding parameters.

The evaluation is conducted using the W\&I + LOCNESS dataset, which provides learner-written sentences with annotated corrections. The model is not fine-tuned; instead, all experiments rely on zero-shot and few-shot prompting.

The study emphasises evaluation rather than deployment. Its goal is not to build a production-ready GEC system, but to analyse how different prompting strategies influence correction behaviour, particularly in relation to over-correction.

The findings should therefore be read as evidence from this particular setup, rather than as a claim about all LLM-based GEC systems.

\section{Overview of the Dissertation Structure}

The remainder of this dissertation is organised as follows.\\

\textbf{Chapter 2} reviews the literature on grammatical error correction (GEC), with particular emphasis on evaluation practices and the limitations of accuracy-based metrics such as ERRANT. It discusses LLM-based GEC, over-correction, pedagogical concerns for second-language (L2) learning, and prompt engineering as a way of influencing correction behaviour. The chapter ends by identifying the gaps that motivate the experiments in this project.\\

\textbf{Chapter 3} turns these gaps into the experimental design used in the project. It defines the research questions, dataset selection, and prompting strategies evaluated in this study, including baseline, instruction-constrained, role-based, and few-shot approaches. The chapter explains how the prompt conditions are compared under a controlled setup, using ERRANT-based accuracy metrics together with change-based measures of textual modification. It also introduces OCI as a proxy for over-correction risk while accounting for fluency and meaning preservation.\\

\textbf{Chapter 4} describes the design and implementation of the experimental pipeline. It explains how original learner sentences are processed, how prompts are applied consistently across conditions, and how outputs are stored and evaluated. The chapter covers both accuracy-based and change-based evaluation components, including the calculation of OCI. It also outlines the analysis procedures used in the study, including prompt refinement, strategy comparison, sentence-length analysis, robustness testing, trade-off analysis, and qualitative inspection, as well as the steps taken to support reproducibility and validity.\\

\textbf{Chapter 5} presents the results and discussion of the experimental study. The analysis is structured around six experiments: prompt refinement, comparison of prompting strategies, sentence-length effects, robustness of the OCI metric, trade-off analysis between correction performance and over-correction, and qualitative analysis of model behaviour. The chapter compares the baseline, instruction, role, and few-shot prompts using ERRANT scores, OCI, edit distance, sentence length, and selected examples.\\

\textbf{Chapter 6} concludes the dissertation by summarising the key findings and contributions of the study. It discusses the implications of prompt engineering for controlling correction behaviour in LLM-based GEC, particularly in educational contexts where minimal and transparent feedback is important. The chapter also outlines the limitations of the study, including dataset and model constraints, and proposes directions for future research, such as adaptive prompting strategies and more fine-grained evaluation of over-correction.\\
